%=====================================================================
\chapter{The Sixteen-Week Schedule}\label{app:schedule}
%=====================================================================
\normalfigures

This appendix is the \textbf{source of record} for the teaching schedule. The
one-page \texttt{AoA\_Fall2026\_Lecture\_Schedule.pdf} distributed to students
is generated from it, and the slide pipeline parses the table below. Change
this first; regenerate the one-page version afterwards.

Twenty-three chapters are taught across sixteen weeks, so seven weeks carry two
chapters. Those pairings are not arbitrary: in every case the second chapter is
the one that cannot be stated honestly without the first --- $O$ and $\Omega$
with $\Theta$, merge sort with the divide-and-conquer pattern that produces it,
graphs with the representation that decides their cost.

\section{The Schedule}

% MACHINE-READ. The slide builder takes the FIRST longtable in this file and
% requires: five columns; a bare integer week; chapters as comma-separated
% digits or N--M ranges; literal --- for an empty assessment cell; and no
% chapter appearing in two weeks.
\rowcolors{2}{white}{lightbg}
% column 3 holds the bold head "Chapters", which needs 1.8 cm at 10.95 pt
\begin{longtable}{@{}L{1.1cm}L{3.2cm}L{1.8cm}L{4.4cm}L{3.9cm}@{}}
\toprule
\rowcolor{primary!15}
\textbf{Wk} & \textbf{Lecture topic} & \textbf{Chapters} & \textbf{Lab} &
\textbf{Assessment} \\
\midrule
\endhead
1 & What an algorithm is; input size and the three cases & 1, 2 & Growth
measured: the ratio test, and the crossover at $n \approx 65$ & --- \\
2 & Asymptotic notation, formally & 3 & Fit a curve; exclude the exponents
either side & --- \\
3 & Correctness and loop invariants & 4 & A bug that is never wrong and never
finishes & --- \\
4 & Recurrences; the master theorem & 5, 6 & Recursion trees against measured
depth and cost & --- \\
5 & Brute force; divide and conquer & 7, 8 & Where $n^{2}$ beats $n \lg n$, and
by how much & --- \\
6 & Quicksort and randomisation & 9 & The adversarial input, and what
randomising costs & \textbf{Quiz 1 (Ch 1--6)} \\
7 & The comparison lower bound; linear-time sorting & 10 & Counting and radix
sort against \texttt{std::sort} & \textbf{Assignment 1 set} \\
8 & Selection and order statistics & 11 & Midterm, then quickselect against a
full sort & \textbf{Midterm (Ch 1--9)} \\
9 & Heaps, search trees and hashing, analysed & 12, 13 & A hash table degrading
under a poor hash & \textbf{Assignment 1 due} \\
10 & Greedy algorithms and exchange arguments & 14 & Huffman coding on real
text; measure the compression & --- \\
11 & Dynamic programming: elements and classics & 15, 16 & Memoised against
bottom-up, timed and counted & --- \\
12 & Amortised analysis & 17 & A growing array: aggregate cost against
worst-case cost & \textbf{Quiz 2 (Ch 10--16)} \\
13 & Graph representation, traversal and spanning trees & 18, 19 & The
matrix/list crossover as density falls & \textbf{Assignment 2 set;
project proposal} \\
14 & Shortest paths & 20 & Dijkstra with and without a priority queue & --- \\
15 & String matching; P, NP and NP-completeness & 21, 22 & Naive, Rabin--Karp
and KMP on the same corpus & \textbf{Assignment 2 due} \\
16 & The semester project & 23 & Implement, measure, defend & \textbf{Project
due; Final (Ch 1--22)} \\
\bottomrule
\end{longtable}

\section{How the Weeks Map to the HEC Outline}

Every clause of the published HEC 2023 outline is reached, and the Course
Specification in the front matter traces each clause to its chapter. The table
below is the same trace read the other way round --- by week, for a course file
that asks when a topic is delivered rather than where it is written.

\rowcolors{2}{white}{lightbg}
\begin{longtable}{@{}L{4.0cm}L{3.0cm}L{7.4cm}@{}}
\toprule
\rowcolor{primary!15}
\textbf{Outline clause} & \textbf{Weeks} & \textbf{Note} \\
\midrule
\endhead
Introduction; role of algorithms in computing & 1 & Ch~1 \\
Analysis on nature of input and size of input & 1 & Ch~2; CLO-1 and CLO-2 \\
Asymptotic notations; $O$, $\Omega$, $\Theta$, $o$, $\omega$ & 2 & Ch~3;
  CLO-5 \\
Sorting algorithm analysis & 5--8 & Ch~8--11, including the lower bound the
  outline does not name \\
Loop invariants & 3 & Ch~4 \\
Recursion and recurrence relations & 4 & Ch~5, 6 \\
Algorithm design techniques; brute force & 5 & Ch~7 \\
Divide-and-conquer; merge, quick sort & 5, 6 & Ch~8, 9; CLO-6 \\
Greedy approach & 10 & Ch~14 \\
Dynamic programming; elements of DP & 11 & Ch~15, 16 \\
Search trees; heaps & 9 & Ch~12 --- analysed, not reimplemented \\
Hashing & 9 & Ch~13 \\
Graph algorithms; shortest paths; sparse graphs & 13, 14 & Ch~18--20; CLO-7 \\
String matching & 15 & Ch~21; CLO-8 \\
Introduction to complexity classes & 15 & Ch~22; CLO-4 \\
\bottomrule
\end{longtable}

\begin{conceptbox}[Two weeks the outline does not ask for]
\textbf{Week 7 (Chapter~10), the comparison lower bound.} The outline says
``sorting algorithm analysis'' and lists the sorts. It does not ask for the
$\bigOm{n\lg n}$ decision-tree proof. It is taught anyway because it is the
first result in the course about a \emph{problem} rather than an algorithm, and
because without it ``merge sort is optimal'' is an assertion rather than a
theorem.

\medskip
\textbf{Week 12 (Chapter~17), amortised analysis.} Also unnamed. Chapter~13's
claim about hash tables and Chapter~19's union-find both depend on it, and
CLO-3 asks for the time complexity of simple algorithms --- which, for a
sequence of operations on a growing structure, is an amortised question. The
front matter states this justification in full.
\end{conceptbox}

\begin{alertbox}[If the semester is shorter than sixteen weeks]
Cut in this order, and say so in the course file:

\begin{tightnum}
  \item \textbf{Chapter~16} (dynamic programming problems) down to two
        problems rather than five. Chapter~15 carries the examinable content.
  \item \textbf{Chapter~21} (string matching) down to naive and Rabin--Karp,
        leaving KMP's prefix function as reading. CLO-8 asks only that a
        string-matching algorithm be traced or implemented.
  \item \textbf{Chapter~12} to a single lecture, since Data Structures has
        already built these structures and this chapter only re-analyses them.
\end{tightnum}

\smallskip
Do \emph{not} cut Chapters~3, 4, 6, 10 or 15. Each is a prerequisite for
something later, and Chapter~10 is the only place the course proves a lower
bound at all.
\end{alertbox}
